\documentclass[11pt]{article}
\usepackage[T1]{fontenc}
\usepackage[utf8]{inputenc}
\usepackage{amsmath,amssymb,amsthm,mathtools}
\usepackage{geometry,graphicx,booktabs,longtable,tabularx,array}
\usepackage{enumitem,needspace,xcolor,tikz,tcolorbox}
\usetikzlibrary{arrows.meta,positioning,calc,fit,backgrounds,decorations.pathreplacing,matrix,patterns}
\tcbuselibrary{breakable}
\usepackage{caption,hyperref,bookmark,cleveref,microtype}
\usepackage{komorebi-publication}
\renewcommand{\KomorebiTitle}{Komorebi Space}
\renewcommand{\KomorebiSubtitle}{Coherent realizations, obstructions, and descent}
\renewcommand{\KomorebiSeries}{Foundations of the Realization Calculus}
\renewcommand{\KomorebiKind}{Monograph}
\renewcommand{\KomorebiEdition}{Expanded monograph edition 1.6.0 / 12 October 2026}
\renewcommand{\KomorebiShortTitle}{Komorebi Space}
\renewcommand{\KomorebiProfile}{monograph}
\KomorebiSetup
\setcounter{tocdepth}{1}
\setcounter{secnumdepth}{2}
\setlist{nosep,leftmargin=1.7em,topsep=4pt}
\hypersetup{pdfkeywords={Komorebi Space,realization calculus,descent,obstruction,Lean,holonomy,closed range}}
\newtheorem{theorem}{Theorem}[section]
\newtheorem{proposition}[theorem]{Proposition}
\newtheorem{lemma}[theorem]{Lemma}
\newtheorem{corollary}[theorem]{Corollary}
\theoremstyle{definition}
\newtheorem{definition}[theorem]{Definition}
\newtheorem{example}[theorem]{Example}
\theoremstyle{remark}
\newtheorem{remark}[theorem]{Remark}
\numberwithin{equation}{section}
\newcommand{\R}{\mathbb R}
\newcommand{\Z}{\mathbb Z}
\newcommand{\Q}{\mathbb Q}
\newcommand{\N}{\mathbb N}
\newcommand{\K}{\mathbb K}
\newcommand{\V}{\mathcal V}
\newcommand{\W}{\mathcal W}
\newcommand{\id}{\operatorname{id}}
\newcommand{\im}{\operatorname{im}}
\newcommand{\ran}{\operatorname{ran}}
\newcommand{\Hom}{\operatorname{Hom}}
\newcommand{\Ob}{\operatorname{Ob}}
\newcommand{\Ext}{\operatorname{Ext}}
\newcommand{\Match}{\operatorname{Match}}
\newcommand{\Exact}{\operatorname{Exact}}
\newcommand{\dist}{\operatorname{dist}}
\newcommand{\Aut}{\operatorname{Aut}}
\newcommand{\Stab}{\operatorname{Stab}}
\newcommand{\Set}{\mathbf{Set}}
\newcommand{\Cat}{\mathbf{Cat}}
\newcommand{\norm}[1]{\left\lVert#1\right\rVert}
\newcommand{\abs}[1]{\left\lvert#1\right\rvert}
\newcommand{\source}[1]{\par\phantomsection\label{src:#1}{\footnotesize\sffamily\color{KomorebiMuted}Source concordance: \hyperlink{claim:#1}{#1}.}\par\smallskip}
\newcommand{\code}[1]{\texttt{\detokenize{#1}}}
\newcommand{\chap}[1]{\par\Needspace{10\baselineskip}\section{#1}}
\newcommand{\partnote}[2]{\clearpage\part{#1}\noindent\begin{minipage}{.92\linewidth}\large #2\end{minipage}\par\vspace{12pt}}
\newtcolorbox{insight}{breakable,colback=KomorebiSoft!55,colframe=KomorebiPine!50,boxrule=.45pt,arc=1pt,left=9pt,right=9pt,top=7pt,bottom=7pt}
\newtcolorbox{scopebox}{breakable,colback=KomorebiWarm,colframe=KomorebiOchre!55,boxrule=.4pt,arc=1pt,left=9pt,right=9pt,top=7pt,bottom=7pt}
\usepackage{float}
\newcommand{\fixedplate}[3]{\begin{figure}[H]\centering\resizebox{.96\linewidth}{!}{\input{figures/#1}}\caption{#2}\label{#3}\end{figure}}
\newcommand{\plate}[3]{\begin{figure}[!htbp]\centering\resizebox{.96\linewidth}{!}{\input{figures/#1}}\caption{#2}\label{#3}\end{figure}}
\tikzset{box/.style={draw=KomorebiPine!70,rounded corners=2pt,fill=KomorebiSoft!50,align=center,inner sep=8pt,font=\small},arr/.style={-{Latex[length=2.2mm]},draw=KomorebiPine,thick},thinarr/.style={-{Latex[length=1.8mm]},draw=KomorebiMuted},lab/.style={font=\small,align=center},dot/.style={circle,fill=KomorebiPine,inner sep=2.3pt}}
\renewcommand{\arraystretch}{1.18}
\raggedbottom
\widowpenalty=7000\clubpenalty=7000
\setlength{\emergencystretch}{3em}

% Compact only the contents entries, without reducing body type size.
\usepackage{etoolbox}
\makeatletter
\patchcmd{\l@section}{1.0em}{0.6em}{}{}
\makeatother

% Generated by checks/edition_stats.py from the current executed reports.
% Generated from current executed reports. Do not hand edit.
\newcommand{\EditionInexactChecks}{18,728}
\newcommand{\EditionFocusedTargets}{7}
\newcommand{\EditionFocusedDeclarations}{223}
\newcommand{\EditionFocusedNewDeclarations}{97}
\newcommand{\EditionFocusedTheorems}{51}
\newcommand{\EditionFocusedReferences}{30}
\newcommand{\EditionClaims}{80}
\newcommand{\EditionReferences}{211}
\newcommand{\EditionWrittenOnly}{0}
\newcommand{\EditionOriginalChecks}{14,224}
\newcommand{\EditionRecoveredChecks}{2,767}
\newcommand{\EditionRevisionChecks}{17,213}
\newcommand{\EditionExpansionChecks}{96,345}
\newcommand{\EditionDeepeningChecks}{13,817}
\newcommand{\EditionTotalChecks}{197,482}
\newcommand{\EditionRegularizationChecks}{34,388}
\newcommand{\EditionWrittenExtensions}{16}
\newcommand{\EditionFormalModules}{138}
\newcommand{\EditionAuditedDeclarations}{4,745}
\newcommand{\EditionNewTheorems}{54}
\newcommand{\EditionFormalTheorems}{1,507}
\newcommand{\EditionNewConstructions}{17}
\newcommand{\EditionReplays}{1}
\newcommand{\EditionNegativeControls}{7}


% Keep readable half-page figures with their surrounding exposition.
\renewcommand{\topfraction}{.88}
\renewcommand{\bottomfraction}{.80}
\renewcommand{\textfraction}{.10}
\renewcommand{\floatpagefraction}{.78}
\setcounter{topnumber}{2}
\setcounter{bottomnumber}{2}
\setcounter{totalnumber}{3}

% WR evidence is generated from separately scoped accepted extension records.
% Generated source-bound WR evidence labels.
\expandafter\def\csname WRstatusWR01\endcsname{Complete written proof. No new Lean counterpart is assigned.}
\expandafter\def\csname WRstatusWR02\endcsname{Complete written proof. No new Lean counterpart is assigned.}
\expandafter\def\csname WRstatusWR03\endcsname{Complete written proof. No new Lean counterpart is assigned.}
\expandafter\def\csname WRstatusWR04\endcsname{Complete written proof. No new Lean counterpart is assigned.}
\expandafter\def\csname WRstatusWR05\endcsname{Complete written proof. No new Lean counterpart is assigned.}
\expandafter\def\csname WRstatusWR06\endcsname{Focused Lean counterpart verified. Examples and prose correspondence remain separately scoped.}
\expandafter\def\csname WRstatusWR07\endcsname{Focused Lean counterpart verified. Examples and prose correspondence remain separately scoped.}
\expandafter\def\csname WRstatusWR08\endcsname{Complete written proof. No new Lean counterpart is assigned.}
\expandafter\def\csname WRstatusWR09\endcsname{Partial focused Lean coverage. The limiting or additional conclusions retain written proofs.}
\expandafter\def\csname WRstatusWR10\endcsname{Complete written proof. No new Lean counterpart is assigned.}
\expandafter\def\csname WRstatusWR11\endcsname{Complete written proof. No new Lean counterpart is assigned.}
\expandafter\def\csname WRstatusWR12\endcsname{Complete written proof. No new Lean counterpart is assigned.}
\expandafter\def\csname WRstatusWR13\endcsname{Focused Lean counterpart verified. Examples and prose correspondence remain separately scoped.}
\expandafter\def\csname WRstatusWR14\endcsname{Focused Lean counterpart verified. Examples and prose correspondence remain separately scoped.}
\expandafter\def\csname WRstatusWR15\endcsname{Complete written proof. No new Lean counterpart is assigned.}
\expandafter\def\csname WRstatusWR16\endcsname{Partial focused Lean coverage. The limiting or additional conclusions retain written proofs.}

\newcommand{\written}[1]{\par\smallskip{\footnotesize\sffamily\color{KomorebiMuted}Extension \hyperref[app:written]{#1}. \csname WRstatus#1\endcsname}\par\smallskip}

\renewcommand{\KomorebiKind}{Lean and evidence guide}
\renewcommand{\KomorebiSubtitle}{Sources, theorem correspondence, and verification}
\renewcommand{\KomorebiEdition}{Companion edition 1.6.0 / 12 October 2026}
\begin{document}
\KomorebiCover
\section*{What this edition establishes}
The eighty KS entries retain their accepted canonical source correspondence. This edition adds a separately verified three-module regularization extension. Four WR groups have complete focused typed counterparts, two have explicitly partial coverage, and ten remain written-only. Current focused execution, retained cumulative evidence, and finite software tests are reported separately.

The original canonical project joins the foundation, moving-space, and Hilbert holonomy sources. Version 1.2.0 closed KS54 through KS58. Subsequent exposition developed KS59 through KS80 from the same accepted source closure. All 156 canonical formal files remain unchanged. The new entrypoint is \path{extension/KomorebiSpaceRegularization.lean}. It imports the three new modules and the three actual inherited dependency modules required by the protected-kernel bridge. Those seven targets, and only that focused closure, have a new accepted clean replay.

All \EditionClaims\ KS catalogue items have source correspondences, comprising \EditionReferences\ resolved declaration references. A reference count is not a count of distinct theorems or discoveries. Several printed statements assemble related declarations. Illustrative calculations and explanatory sentences are not automatically promoted to separately formalized claims.

\begin{scopebox}
\textbf{Retained formal baseline: version 1.2.0}\par
The canonical project has \EditionFormalModules\ production modules and one root. The completed clean source replay records, imported axiom audit, positive examples, and deliberate rejection controls are in \path{audit/formal/}. The audit covers \EditionAuditedDeclarations\ project declarations, including compiler-generated declarations. Its permitted axiom set is \code{propext}, \code{Classical.choice}, and \code{Quot.sound}. The pinned external dependency cache is reused rather than rebuilt.
\end{scopebox}

\subsection*{The distinctions that remain important}
A successful compiler run certifies the typed declarations, not a translation of every word of the manuscript. A correct theorem can have narrower hypotheses than an informal description. Its application hypotheses still need to hold in a concrete mathematical model. Formal verification also does not determine historical novelty.

Accordingly, this guide separates retained compiler evidence, imported axiom dependencies, theorem-to-paper correspondence, finite implementation tests, and the remaining mathematical boundaries. The actual Lean types remain authoritative for formal scope.

\subsection*{Files to start with}
The historical cumulative entrypoint is \path{formal/KomorebiSpace.lean}. The focused extension entrypoint is \path{extension/KomorebiSpaceRegularization.lean}, with its release summary at \path{audit/extension/release.json}. The canonical sources are in \path{formal/KomorebiSpace/}. The source map is \path{data/claim_ledger.json}, with a CSV companion. The retained formal release summary is \path{audit/formal/release.json}. Each replay has its own complete command receipt, source hashes, axiom report, and logs. The root progress ledger identifies the current restart point.

\clearpage\section{One canonical project, preserved provenance}
\subsection{F: the foundation}
The unchanged archive \path{Komorebi_Space_Lean_v0.20.0.zip} contains 122 production modules and a root. Its historical inventory reports 1,299 named theorem declarations and 696 constructions, including auxiliary results. This tree develops views, partial extension, additive obstruction, tower repair, residual cascades, nonabelian cocycles, effective descent, the composition-site stack, presentation transport, and finite-dimensional amplified holonomy repair.

Its historical strategy included one full cumulative build and a subsequent archive replay reusing a checked 117-module baseline while rebuilding five additions and the root. Those historical records are not relabeled as two clean cumulative replays.

\subsection{M: the focused moving-space closure}
The unchanged archive \path{Komorebi_Moving_Exact_Spaces_v0.32.0.zip} contains 11 production modules and a root. Its historical inventory reports 151 named theorem declarations and 47 constructions. Six modules are inherited dependencies of that focused closure and five are additions. Its recorded verification includes two clean builds of the supplied closure, axiom requests, and negative controls.

The higher version number does not make this archive a replacement for the foundation. It supplies closed-range repair criteria, moving-kernel estimates, constructed perturbed splittings, actual kernel transport, and the quantitative-descent bridge.

\subsection{S: Hilbert holonomy}
The separately recovered \path{HilbertHolonomyRepair.lean} implements amplified holonomy repair with Hilbert fibers and its closed-range necessity criterion. It lay outside both historical archive replay scopes. Version 1.2.0 included it in the canonical source closure and compiled it with its actual dependencies. That accepted evidence is retained here and bound to unchanged source bytes, not borrowed from an unrelated historical report.

\subsection{I: canonical sources and the five version 1.2.0 additions}
The five modules added in version 1.2.0 are \path{FinitePrefixBoundary.lean}, \path{DyadicSupport.lean}, \path{ProtectedHilbertRepair.lean}, \path{ProtectedLinearCertificates.lean}, and \path{FiniteRealizationExistence.lean}. They contain \EditionNewTheorems\ explicit named theorems and lemmas and \EditionNewConstructions\ indexed constructions. These counts include supporting results and do not assert that the classical principles used are newly discovered.

\begin{center}
\begin{tabularx}{\linewidth}{@{}p{.58\linewidth}r@{}}
\toprule Source-selection accounting & Modules\\\midrule
Foundation production sources & 122\\
Focused production sources & 11\\
Separate Hilbert holonomy source & 1\\
Identical shared module counted once & $-1$\\
Version 1.2.0 formalization modules & 5\\\midrule
Canonical production modules & \EditionFormalModules\\
Canonical root & 1\\\bottomrule
\end{tabularx}
\end{center}
The shared module is \code{KernelProjectionBounds}. Its two copies are byte-identical. No inherited production source needed a mathematical or elaboration repair to compile in the canonical union. Both original roots and all frozen archives are preserved separately.

\section{The five formalization gaps that were closed}
\subsection{KS54: finite prefixes do not supply a global point}
The formal model uses the existing \code{Resolution} with fibers \(X_n=\{m\in\N:n\le m\}\) and inclusion restrictions. For every \(N\), the actual partial-section type on the prefix \(0,\ldots,N\) is inhabited by the constant value \(N\). A global point would have one integer value at every coordinate, contradicting the membership condition at its successor.

The proof additionally transfers emptiness to the existing ambient-section type and rejects the existence of a compatible matching family of all prefix sections. This makes the distinction precise: separate prefix witnesses are not matching prefix data.

\subsection{KS55: the dyadic support criterion}
The induction is carried out on the actual \code{DyadicResidual.System}, retaining its changing operator fibers and the starting scale. It proves
\[
 2^N\nmid(k+1)\quad\Longrightarrow\quad (Q_Nr)_k=0.
\]
The corresponding reconstructed coordinate is exact. A nonzero remainder coordinate must lie on the allowed congruence support. The additional valuation form gives exact recovery when \(N>\nu_2(k+1)\).

The induction uses the actual parity-lift identity. It does not replace the varying-fiber system by a scalar model. Neither support theorem asserts that every allowed coordinate is occupied or that coordinatewise recovery implies uniform-norm convergence.

\subsection{KS56: nearest repair with a protected observable}
The source is a complete real Hilbert space. Both the defect \(D\) and protected observable \(p\) are continuous linear maps into the declared normed spaces. Corrections belong to the actual closed kernel \(M=\ker p\), and the operator being studied is \(B=D|_M\).

Feasibility is proved equivalent to the existence of an exact answer preserving \(p(x)\). For feasible \(x\), a raw permitted correction is projected off \(\ker B\) inside \(M\). The resulting displacement is the unique orthogonal permitted correction with image \(Dx\). The actual repaired answer satisfies exactness, preservation, nearest-point uniqueness, and the Pythagorean identity.

No closed-range assumption is needed for nearest-point existence. With a Banach residual target, a uniform residual bound on all feasible inputs is equivalent to closedness of \(\operatorname{ran}B\). The inverse norm is proved to be the least admissible nonnegative constant, including the zero-operator case. Closedness of \(D(E)\) alone is not substituted for this restricted image condition.

\subsection{KS57: literal field-linear certificates}
For linear maps over an arbitrary field, the generic theorem supplies either a protected exact answer or normalized algebraic functionals satisfying the obstruction equations. Algebraic dual separation proves the negative alternative without assuming an elimination routine is correct.

The finite-coordinate specialization proves the literal row-vector statement
\[
 \lambda A=\mu P,\qquad \lambda Ax=1.
\]
It also proves exclusivity and exhaustiveness, identifies the answer fiber as an affine translate of \(\ker A\cap\ker P\), and proves its dimension by rank-nullity for the literal vertically stacked matrix. Zero-dimensional index sets are permitted.

The printed proof uses row reduction, whereas the formal proof uses dual separation. Both establish the stated mathematical alternative. Neither is a universal refinement theorem for the separate Python solver, and neither replaces an integral feasibility problem by a rational one without changing the question.

\subsection{KS58: finite fibers and persistent images}
An explicit functor from \(\N^{op}\) is constructed from the original resolution maps. A coordinate-preserving equivalence identifies its sections with the original \code{Resolution.Point}. The pinned mathlib finite cofiltered-system theorem proves nonemptiness when every fiber is finite and nonempty.

The eventual-range machinery then identifies the range of evaluation at level \(n\) with the intersection of all later restriction images. This proves the persistent-image extension criterion, not surjectivity of evaluation onto every fiber. The printed sequential stabilization proof and the formal cofiltered-system proof have different routes but the same declared conclusion.

\section{Replays, axiom reports, and controls}
\subsection{Pinned runtime of the retained baseline}
For the accepted version 1.2.0 replay, the portable runtime was restored from seven retained offline archive parts. That restoration is recorded evidence, not a new operation in the present editorial expansion. Their payload hashes and CRCs were checked, as was the assembled archive stream. The restored Lean version is 4.34.0-rc2 at commit
\begin{center}\small\path{6a10ac8c22beadecabdbb0919c2b50214762f91d}.\end{center}
The mathlib revision is
\begin{center}\small\path{aeb72f1e9981b6d35b426b16dca7759be3e7ba06}.\end{center}
The verifier checks all nine dependency revisions against the delivered pin manifest. No network operation or floating dependency update is required by the offline replay.

\subsection{Accepted baseline source compilation}
Each completed clean replay begins with an empty local build directory. It resolves the complete local import graph, rejects missing modules and cycles, checks that the root reaches every production file, and compiles each source with one worker, warnings treated as errors, and automatic implicit variables disabled. Every command has an exit code, timing, log hash, and source-bound output record.

The release summary identifies the completed clean replay records. The sources are checked for mutation across each run. A timed-out or failed command is not a pass. Failed development attempts are preserved separately rather than used as final compiler evidence.

\subsection{Imported axiom audit}
The audit enumerates declarations by their defining project module, including generated declarations. It asks Lean for transitive axiom dependencies and rejects any project declaration using an axiom outside \code{propext}, \code{Classical.choice}, and \code{Quot.sound}. Unsafe declarations are rejected as well. The retained environment inventory contains \EditionAuditedDeclarations\ audited project declarations. The lexical source inventory separately counts \EditionFormalTheorems\ explicit theorem and lemma declarations.

The dependency query uses Lean's imported axiom machinery. The compiler, kernel, and imported cached mathematical dependencies remain in the trust boundary. This is not an independent kernel implementation or a full source rebuild of mathlib. Source hashes authenticate the bytes recorded by this package, not an external execution authority.

\subsection{Positive and negative controls}
The positive Lean regression file instantiates protected rational answers, a normalized obstruction, distinct exact answers without the protected constraint, infeasibility under the identity observable, and finite Boolean resolution examples. It also combines separate prefix inhabitance with the proved absence of a global point.

The final control suite has \EditionNegativeControls\ rejection cases. False arithmetic, a proof placeholder, the false finite-prefix implication, universal protected feasibility, an unnormalized certificate, and a wrong preserved value must be rejected by the compiler. A deliberately introduced custom axiom is legal Lean syntax but must be rejected by the dependency policy. That last result is an audit rejection, not a claim that Lean forbids declaring axioms.

\subsection{Exact finite implementation checks}
The independent publication checker executes \EditionOriginalChecks\ assertions. The recovered moving-space checker executes \EditionRecoveredChecks. The revision suite executes \EditionRevisionChecks, including 969 rational certificate problems, independent minor-based rank checks, varying matrix-fiber dyadic examples, and nine implementation rejection controls. The expansion suite adds \EditionExpansionChecks\ exact assertions covering prescribed branches, finite observation, repair mismatch, loop reconstruction, noncommutative boundary comparisons, and uncertainty bounds. The deepening suite adds \EditionDeepeningChecks\ exact assertions for variable-rate paths, closure-boundary approximations, directed costs, and the finite phase extension. The regularization suite adds \EditionRegularizationChecks\ assertions. The inexact-certificate suite adds \EditionInexactChecks\ assertions. Together the seven rerun suites contain \EditionTotalChecks\ exact finite assertions. The expansion suite also verifies that all formal files retain their baseline hashes.

These computations check their stated finite models and implementations. The general results identified by principal KS entries come from the formal declarations under their hypotheses. The additional infinite-dimensional example is proved in the manuscript. Neither type of universal conclusion is obtained by extrapolating finite samples. The Python Gaussian proposer remains tested software rather than a universally verified Lean program.

\section{The preceding exposition expansion}
The twelve preceding exposition entries KS59--KS70 did not enlarge the formal source inventory. They expose bridges that were already implemented but insufficiently developed in the reader edition.

\begin{center}
\begin{tabularx}{\linewidth}{@{}p{.18\linewidth}X@{}}
\toprule Entries & Added mathematical development\\\midrule
KS59--KS60 & Stronger prescriptions, quotient naturality, shrinking ambiguity, and a three-view integer choice problem.\\
KS61--KS64 & Scheduled transformations, eventual equality, faithfulness from surjective point projections, the finite-observation image theorem, and cofinal transport.\\
KS65--KS67 & Split exact/residual coordinates, the off-diagonal repair mismatch, its composition law, and bounded-gap residual sampling.\\
KS68 & A concrete failure of literal full-framing invariance, contrasted with marked boundary identifications.\\
KS69--KS70 & Measured residuals, operator uncertainty, and the strict-radius repair theorem when the old kernel is preserved.\\
\bottomrule
\end{tabularx}
\end{center}

The two loop chapters are fully calculated worked applications. The first distinguishes changing inconsistent edge data from changing a vertex candidate. The second distinguishes root-based reconstruction from nearest repair in an explicitly declared Euclidean norm. Their finite calculations are tested independently, but the edition does not assign a new Lean theorem to every numerical specialization or exercise.

The source map has \EditionClaims\ principal items and \EditionReferences\ resolved references. All new references occur in the retained axiom inventory and have unchanged source hashes. Printed proofs remain an editorial mathematical layer. Their presence in the monograph is not a claim of automatic prose-to-Lean translation.

\section{The retained convergence and obstruction development}
KS71--KS75 use the retained \path{Exactification}, \path{LinearExactification}, \path{Profiles}, and \path{RepairExamples} modules. The key logical separation is between residual decay, control of total movement, and a comparison between the limits of different paths. The printed examples instantiate these distinctions with exact scalar coordinates. They are not newly compiled instantiations.

KS76--KS77 use \path{ObstructionFloor}. In a normed additive group, zero distance to the correction subgroup means membership in its closure. The manuscript's diagonal Hilbert-space model proves directly that this closure can be strictly larger than the actual correction image. Its infinite-dimensional nonexistence and escape arguments are written proofs, not conclusions inferred from the finite checker and not a new Lean instance.

KS78--KS79 use \path{WeightedGeometry}. They establish extended directed-distance laws only after an independent arrow cost has been declared. The source's zero-distance setoid does not identify categorical isomorphism classes. The free-path example supplies a concrete printed counterexample to that identification.

KS80 uses \path{ProjectiveExample}, with an additive phase in \(\mathbf F_2\). The actual retained theorem concerns the cocycle law and the impossibility of a scalar coboundary. Its matrix realization and eight-element extension are fully worked and exhaustively tested finite interpretations. No new representation-theory module or arbitrary higher-cohomology theorem is claimed.

The model-to-evidence ledger is \path{data/deepening_models.json}. It distinguishes the general retained declarations from printed instantiations, finite arithmetic tests, and bounds on an infinite tail. None of these additions alters the accepted formal tree.

\section{A separately verified regularization extension}
The focused extension is not a replacement for the canonical library. Its clean replay compiles three inherited modules, three new modules, and one root. The inherited files are \path{KernelProjectionBounds.lean}, \path{HilbertKernelProjection.lean}, and \path{ProtectedHilbertRepair.lean}. Their original source hashes are checked. No broad import of an unrelated project is used to hide dependencies.

\subsection{Constructing the optimizer}
\path{RegularizedHilbertRepair.lean} defines the actual normal operator
\[
 N=T^*T+\alpha I.
\]
The associated continuous bilinear form is proved coercive for $\alpha>0$. The pinned Lax--Milgram theorem constructs a continuous linear equivalence, whose inverse defines $a_\alpha$. Existence is derived, not stored as an assumed field. The module proves the normal equation, exact energy identity, global minimality and uniqueness, parameter monotonicity, fixed-budget optimality, and the normal-residual and data-noise inequalities.

The manuscript proves invertibility using self-adjointness and a closed-range argument. The formal proof uses the equivalent coercive-form route. This is a documented proof-route difference, not a different minimizer or an omitted existence obligation.

\subsection{Certifying a proposal and its dual witness}
\path{RegularizationCertificates.lean} defines the unconstrained quadratic dual objective and proves the exact two-square gap identity. It derives weak duality, attainment at the residual of the constructed minimizer, a correction-error certificate, and the residual-witness specialization. The same actual optimizer supplies a normalized sharp support witness at its own correction budget. This last result does not formalize the more general budget-attainment and weak-compactness arguments WR01--WR05.

\path{ProtectedRegularizedRepair.lean} instantiates the construction on the existing \path{ProtectedHilbertRepair.restricted} operator, whose source is the actual closed kernel of the protected observable. The bridge proves that the restricted adjoint is the projected ambient adjoint. Every proposed correction in that subtype preserves the observable exactly. Measured residual data are retained independently of $Df$, so the actual repaired residual is not silently equated with a measured fit.

\subsection{The precise WR coverage}
\begin{center}
\begin{tabularx}{\linewidth}{@{}p{.17\linewidth}p{.22\linewidth}X@{}}
\toprule Groups & Coverage & Boundary\\\midrule
WR06--WR07 & Complete focused counterpart & Constructed optimizer, energy, uniqueness, parameter and actual-budget results. Coordinate examples remain separate.\\[4pt]
WR13--WR14 & Complete focused counterpart & Universal solver inequalities, dual gap and protected bridge. Scalar sharpness examples are printed and tested.\\[4pt]
WR09, WR16 & Partial focused counterpart & Pointwise data and solver inequalities only. The limiting conclusions still use written WR08.\\[4pt]
WR01--WR05, WR08, WR10--WR12, WR15 & Written proofs only & No new typed coverage of compactness, the full budget profile, infinite limits, diagonal asymptotics, or coordinate-basis invariance.\\\bottomrule
\end{tabularx}
\end{center}
The metadata resolver checks \EditionFocusedReferences\ exact references against the freshly imported declaration types and source hashes. That is an editorial correspondence audit, not a machine-certified translation of natural language. The full list appears in \path{audit/regularization/written_scope.json}.

\subsection{Executed focused verification}
The accepted focused replay has \EditionFocusedTargets\ successful source targets. It audits \EditionFocusedDeclarations\ declarations originating in the six supplied production modules, including \EditionFocusedNewDeclarations\ originating in the three new modules. Generated declarations are included. The new source inventory contains \EditionFocusedTheorems\ explicit named theorems and lemmas, including auxiliary results. These counts are not added blindly to the older cumulative inventory because the focused audit also contains inherited declarations.

The compiler is the restored pinned Lean 4.34.0-rc2 runtime. All nine dependency revisions are checked. Every focused project module is recompiled from an empty local output directory, with warnings treated as errors. Imported mathlib caches remain trusted and are not rebuilt. All imported project declarations are inspected for unapproved axioms and unsafe definitions. Only \code{propext}, \code{Classical.choice}, and \code{Quot.sound} occur.

A positive regression file exercises construction, preservation, solver bounds, and duality. Four compiler controls reject zero-parameter coercivity, omission of the parameter gain, reversed duality, and a proof placeholder. A separate policy control rejects a deliberately introduced custom axiom. The latter is an audit failure by design, not an assertion that Lean syntax forbids axioms.

The attempted cumulative run \path{audit/formal/baseline16/} was explicitly stopped before completion. Its status remains INTERRUPTED and it is not accepted as pass evidence. The focused run was performed separately. The review report retains development failures and any superseded focused attempts.

\subsection{Exact software, without a circular optimizer check}
The optimizer example remains \path{examples/budgeted_repair.py}. The new \path{examples/inexact_repair_certificate.py} accepts a proposed ambient correction and an optional independent dual witness. It checks protection and a full kernel basis before returning the actual candidate, fit residual, model residual, normal covector, Gram-corrected Hilbert residual, duality gap, and squared error bounds. It does not invoke an exact optimizer to certify the proposal.

A separate oracle does solve small full ambient constrained systems, using Leibniz determinants and Cramer formulas rather than the program's nullspace elimination route. The new finite suite covers 1,458 candidate comparisons, 909 independent dual witnesses, 286 nonorthogonal basis changes, 144 scalar solver/data cases, and eighteen structured three-dimensional, four-dimensional, and redundant-prescription cases. Ten false-claim controls, twelve invalid-input cases, and four tampered-output cases are rejected. These are exact finite tests, not universal refinement or a complexity theorem.

\section{Reading theorem-to-paper correspondence}
The stable KS identifiers organize exposition rather than declaration counts. The source resolver matches names, defining modules, and hashes against the retained accepted compiled-environment inventory. Old F, M, and S references retain their original source paths and are also matched to byte-identical canonical copies. I references point directly to canonical sources, including both inherited modules and the five version 1.2.0 additions.

The consequential hypotheses are retained: transition versus composition covers, feasibility versus uniqueness, identified boundaries, closed-range stability, and the residual bounds required for perturbation. These distinctions remain explicit in the monograph and in the actual Lean types.

\clearpage\section{The correspondence catalogue}
The version 1.2.0 formalization review checked the five written additions against their actual types, including the restricted kernel in KS56, normalization and coefficient field in KS57, and finite nonempty fibers in KS58. The inherited correspondence is preserved. The expanded correspondence is checked against actual declarations and unchanged source hashes in the same accepted baseline. It does not claim a new compiler run. No automatic semantic equivalence checker for natural-language proofs is claimed.

Each entry records the editorial role, actual declaration, source line, and short source-hash prefix. The JSON ledger supplies complete paths and hashes. Repeated references are intentional where a source supports several parts of the exposition.

\textbf{F} is foundation provenance. \textbf{M} is focused moving-space provenance. \textbf{S} is the recovered Hilbert source. \textbf{I} is a canonical source module. All referenced source files belong to the accepted version 1.2.0 compiled closure.
\input{guide_entries}

\clearpage\section{Reproduction and handoff}
\subsection{Verify the delivered payload first}
From the extracted package root, run
\begin{quote}\small\ttfamily
python3 checks/verify\_payload.py
\end{quote}
This checks the manifest and standalone-font exclusion without rewriting the package. It is an integrity check, not a Lean proof check.

\subsection{A new formal replay}
Supply the pinned portable runtime and choose a new label:
\begin{quote}\small\ttfamily
python3 formal/scripts/verify.py \textbackslash\par
\hspace*{1em}--runtime /path/to/portable-runtime \textbackslash\par
\hspace*{1em}--label independent1
\end{quote}
The verifier refuses to overwrite an existing audit or local build directory. Its command log records the result of each compilation, axiom query, regression file, and rejection control. Local compiled objects are disposable and are not required in the source distribution. The runtime and external dependency cache are not embedded in this package.

The \path{lakefile.toml} and \path{lean-toolchain} describe the same project for a conventional Lean setup. The fully recorded replay path is the supplied offline verifier, not an unexecuted claim about every installation method.

\subsection{Replay the focused extension}
For the separately delimited regularization closure, run
\begin{quote}\small\ttfamily
python3 extension/scripts/verify.py \textbackslash\par
\hspace*{1em}--runtime /path/to/portable-runtime \textbackslash\par
\hspace*{1em}--label independent-focused1
\end{quote}
This verifier refuses to overwrite a label. Its receipt lists the actual source closure, dependency pins, source hashes, all compiler commands, axiom report and controls. The project-local \path{extension/.work/} directory is disposable and excluded from the source release. A newly produced PASS receipt must be deliberately bound in the focused release summary before publication metadata can consume it.

\subsection{Regenerate the reader set}
The root \path{build.sh} reruns source integrity, correspondence, finite checks, figure generation, evidence counts, and all three PDF builds. It does not silently execute another Lean replay. It consumes the completed source-bound formal release record. Re-render the resulting PDFs using \path{checks/render_review.py} and inspect the pages before renewing the hash-bound visual-review record.

Rebuilding may update current timestamps and files. Historical archives and replay evidence retain their own dates and scope. A modified source needs a new formal replay and release binding before it can inherit current verification status.

\section{Remaining scope}
Version 1.2.0 closed the identified canonical source-integration and five-item formalization tasks. Version 1.6.0 preserves that baseline and adds a newly verified focused regularization closure. The WR coverage table above is deliberately nonuniform. Weak compactness, infinite regularization limits, the diagonal budget asymptotic, and coordinate-covector invariance still have written rather than new typed proofs. Universal-view doctrines, higher stacks, unrestricted semantic realization, nonlinear repair, and general arithmetic correspondences remain separate questions. No new full cumulative replay, independent kernel, full source rebuild of mathlib, native Windows execution, or universal Python refinement is claimed.

The progress ledger records the expansion, failures, checks, and exact restart point. The publication version describes this reader edition, not completion of every proposed Komorebi extension.
\end{document}
